\documentclass[a4paper, serif, 10pt]{moderncv}

%% moderncv
% color and style
\moderncvcolor{black}
\moderncvstyle{banking}

%% page layout/margins
\usepackage[top=2.5cm, bottom=2.5cm, left=1.5cm, right=1.5cm]{geometry}

\nopagenumbers{}

%% Babel config for Norwegian language
% ref:
% - https://latex3.github.io/babel/guides/locale-norwegian.html
\usepackage[norsk]{babel}

%% fontspec
\usepackage{fontspec}

\IfFontExistsTF{Linux Libertine}{
  \setmainfont[Ligatures={TeX, Common}, Numbers=OldStyle]{Linux Libertine}
}{}
\IfFontExistsTF{Linux Libertine O}{
  \setmainfont[Ligatures={TeX, Common}, Numbers=OldStyle]{Linux Libertine O}
}{}

%% CTeX
% CJK support, used to show CJK characters in the resume
%
% - heading=false: disable ctex heading
% - punct=kaiming: use kaiming punctuations styles for CJK
% - scheme=plain: use plain scheme, do not override \normalsize font size
% - space=auto: space settings for CJK characters
%
% CJK fonts are resolved per locale language, e.g. Noto Serif/Sans CJK JP
% for Japanese resumes, with fallback to the next region if not installed
%
% ref:
% - http://ctan.mirrorcatalogs.com/language/chinese/ctex/ctex.pdf
\usepackage[UTF8, heading=false, punct=kaiming, scheme=plain, space=auto]{ctex}

\IfFontExistsTF{Noto Serif CJK SC}{
  \setCJKmainfont{Noto Serif CJK SC}
}{}
\IfFontExistsTF{Noto Sans CJK SC}{
  \setCJKsansfont{Noto Sans CJK SC}
}{}

%% line spacing
\usepackage{setspace}
\setstretch{1.125}

%% URL styling - use normal text instead of monospace
\urlstyle{same}

%% Patch \httplink and \httpslink to handle full URLs
% The original moderncv macros always prepend http:// or https:// to URLs.
% This causes issues when the URL already has a protocol (e.g., https://example.com)
% resulting in malformed URLs like https://https://example.com.
% This patch checks if the URL already contains "://" and uses it directly if so.
% Note: We use \str_set:Nx (x-expansion) to expand macros like \@homepage first.
\ExplSyntaxOn
\str_new:N \g_yamlresume_url_str

\RenewDocumentCommand{\httpslink}{O{}m}{
  \str_set:Nx \g_yamlresume_url_str {#2}
  \str_if_in:NnTF \g_yamlresume_url_str {://}
    { \url{#2} }
    { \url{https://#2} }
}

\RenewDocumentCommand{\httplink}{O{}m}{
  \str_set:Nx \g_yamlresume_url_str {#2}
  \str_if_in:NnTF \g_yamlresume_url_str {://}
    { \url{#2} }
    { \url{http://#2} }
}
\ExplSyntaxOff

\name{Ingrid Solberg}{}
\title{Maskiningeniør med fokus på produktutvikling og bærekraft}
\phone[mobile]{+47 987 65 432}
\email{ingrid.solberg@example.no}
\homepage{https://www.ingridsolberg.no}

\address{Storgata 42, Oslo, Oslo, Norge, 0182}{}{}

\extrainfo{{\small \faLinkedin}\ \href{https://www.linkedin.com/in/ingridsolberg}{@ingridsolberg} {} {} {} • {} {} {} 
{\small \faGithub}\ \href{https://github.com/ingridsolberg}{@ingridsolberg}}

\begin{document}

\maketitle

\section{Grunnleggende informasjon}

\cvline{}{\begin{itemize}
\item Maskiningeniør med mastergrad fra NTNU og fem års erfaring fra produktutvikling, konstruksjon og teknisk prosjektledelse
\item Erfaren med CAD, FEM-analyse og materialvalg for komponenter i krevende industrielle miljøer
\item Dokumentert evne til å lede tverrfaglige team fra konsept til produksjon, med vekt på kvalitet, kostnad og bærekraft
\item Sterk interesse for energieffektivisering, sirkulærøkonomi og digitale verktøy i ingeniørfaglig arbeid
\end{itemize}}
\section{Utdanning}

\cventry{aug. 2017–juni 2022}
        {Master, Maskinteknikk, Poeng: A+}
        {Norges teknisk-naturvitenskapelige universitet}
        {\url{https://www.ntnu.no}}
        {}
        {\begin{itemize}
\item Fordypning i produktutvikling og beregningsmekanikk
\item Masteroppgave om lette konstruksjoner i aluminium for offshore-applikasjoner
\item Gjennomførte tverrfaglige prosjekter med fokus på bærekraftige materialvalg
\end{itemize}
\textbf{Kurs}: Konstruksjonsteknikk, Fasthetslære, Strømningslære, Termodynamikk, Materialteknologi, Produktutvikling, FEM-analyse, Prosjektledelse}
\section{Arbeidserfaring}

\cventry{jan. 2023–Nå}
        {Senior maskiningeniør}
        {Hydro Extrusion Norge AS}
        {\url{https://www.hydro.com/no-NO/}}
        {}
        {\begin{itemize}
\item Leder teknisk utvikling av ekstruderte aluminiumsprofiler for transport- og byggsektoren
\item Ansvarlig for konstruksjonsberegninger, toleranseanalyser og validering av nye produkter
\item Koordinerer tverrfaglige prosjekter mellom produksjon, kvalitet og leverandører
\item Driver kontinuerlig forbedring av produksjonsprosesser med vekt på energibesparelse
\end{itemize}
\textbf{Nøkkelord}: Aluminium, Produktutvikling, FEM, Kvalitetssikring, Lean}

\cventry{aug. 2020–des. 2022}
        {Maskiningeniør}
        {Aker Solutions AS}
        {\url{https://www.akersolutions.com}}
        {}
        {\begin{itemize}
\item Konstruerte og dimensjonerte mekaniske komponenter for undervannsinstallasjoner
\item Utførte FEM- og utmattingsanalyser i samsvar med NORSOK-standarder
\item Utarbeidet tekniske tegninger, spesifikasjoner og endringsnotater
\item Deltok i verifikasjon, testing og oppfølging av leverandører
\end{itemize}
\textbf{Nøkkelord}: Undervannsteknologi, FEM-analyse, NORSOK, SolidWorks, Teknisk dokumentasjon}

\cventry{juni 2019–aug. 2019}
        {Prosjektingeniør (sommerjobb)}
        {SINTEF Energi AS}
        {\url{https://www.sintef.no/energi}}
        {}
        {\begin{itemize}
\item Analyserte energibruk i industrielle prosesser og foreslo forbedringstiltak
\item Utarbeidet simuleringsmodeller for varmegjenvinning i pilotanlegg
\item Presenterte resultater for prosjektteam og eksterne samarbeidspartnere
\end{itemize}
\textbf{Nøkkelord}: Energiøkonomi, Varmegjenvinning, Simulering}
\section{Språk}

\cvline{Norsk}{Morsmål eller tospråklig nivå \hfill \textbf{Nøkkelord}: Morsmål, Bokmål}
\cvline{Engelsk}{Fullt profesjonelt nivå \hfill \textbf{Nøkkelord}: C1, Teknisk engelsk}
\cvline{Tysk}{Grunnleggende nivå \hfill \textbf{Nøkkelord}: A2}
\section{Ferdigheter}

\cvline{Mekanisk konstruksjon}{Ekspert \hfill \textbf{Nøkkelord}: SolidWorks, Siemens NX, AutoCAD, Teknisk tegning, Toleranseanalyse}
\cvline{Beregning og analyse}{Avansert \hfill \textbf{Nøkkelord}: ANSYS Mechanical, Abaqus, FEM-analyse, Utmattingsanalyse, CFD}
\cvline{Materialer og produksjon}{Avansert \hfill \textbf{Nøkkelord}: Aluminium, Stål, Kompositter, Additiv tilvirkning, Lean produksjon}
\cvline{Prosjektledelse}{Avansert \hfill \textbf{Nøkkelord}: Tverrfaglig samarbeid, Teknisk dokumentasjon, Kvalitetssikring, Risikostyring, Bærekraftig design}
\section{Utmerkelser}

\cventry{juni 2022}
        {NTNU}
        {Beste masteroppgave i maskinteknikk}
        {}
        {}
        {Tildelt for masteroppgaven om lette aluminiumskonstruksjoner for offshore-applikasjoner.}
\section{Sertifikater}

\cventry{mars 2021}
        {Dassault Systèmes}
        {Certified SolidWorks Professional}
        {\url{https://www.solidworks.com/certification}}
        {}
        {}

\cventry{nov. 2023}
        {Project Management Institute}
        {PMP – Project Management Professional}
        {\url{https://www.pmi.org/certifications/project-management-pmp}}
        {}
        {}
\section{Publikasjoner}

\cventry{sep. 2021}
        {Energieffektivisering av varmepumper i norske næringsbygg}
        {Norsk Energi}
        {\url{https://www.norskenergi.no/}}
        {}
        {\begin{itemize}
\item Undersøker hvordan varmepumper kan redusere energibehovet i næringsbygg
\item Presenterer beregningsmodeller og anbefalte tiltak for eksisterende bygningsmasse
\item Diskuterer lønnsomhet og bærekraftige installasjonsløsninger
\end{itemize}}
\section{Prosjekter}

\cventry{jan. 2021–juni 2022}
        {Utvikling av en lettvektsramme i aluminium for offshore-applikasjoner.}
        {Bærekraftig lettvektramme for offshore}
        {\url{https://www.ntnu.no/prosjekter/lettvektramme}}
        {}
        {\begin{itemize}
\item Designet og optimaliserte rammekonstruksjonen med hensyn til styrke og vekt
\item Gjennomførte FEM-analyse og utmattingsvurdering av kritiske knutepunkter
\item Validerte konstruksjonen gjennom laboratorietester og dokumenterte resultatene
\end{itemize}
\textbf{Nøkkelord}: Lettvektsdesign, Aluminium, FEM-analyse, Offshore, Bærekraft}

\cventry{juni 2019–aug. 2019}
        {Analyse og forbedring av varmegjenvinning i et industrielt pilotanlegg.}
        {Energioppgradering av pilotanlegg}
        {\url{https://www.sintef.no/prosjekter/energioppgradering}}
        {}
        {\begin{itemize}
\item Kartla energistrømmer og identifiserte tapskilder i prosessen
\item Modellerte varmegjenvinning og foreslo tiltak med god lønnsomhet
\item Bidro til teknisk rapport og anbefalinger for videre drift
\end{itemize}
\textbf{Nøkkelord}: Varmegjenvinning, Energianalyse, Simulering}
\section{Interesser}

\cvline{Friluftsliv}{Fjellturer, Langrenn, Kajakk}
\cvline{Teknologi}{3D-printing, Arduino, Bærekraftig energi}
\cvline{Musikk}{Gitar, Kor}
\section{Frivillig arbeid}

\cventry{sep. 2021–Nå}
        {Fagansvarlig for bærekraft}
        {Tekna – Teknisk-naturvitenskapelig forening}
        {\url{https://www.tekna.no}}
        {}
        {\begin{itemize}
\item Arrangerer faglige møter og workshops om bærekraftig ingeniørpraksis
\item Bidrar med faglige innspill til studenter og nyutdannede ingeniører
\item Koordinerer samarbeid mellom lokale bedrifter og foreningen
\end{itemize}}
    
\end{document}